\chapter{Integration of measures}

Throughout this chapter, T denotes a locally compact space, \(\mu\) a positive measure on T. For every subset A of a set E, \(\varphi_{A}\) denotes the characteristic function of A (if no confusion can result thereby). By numerical function, we always mean a function taking its values in \(\overline{R}\) , thus possibly taking on the values \(+\infty\) and \(-\infty\) . The set of positive numerical functions defined on E will be denoted \(\mathcal{F}_{+}(E)\) , or simply by \(F_{+}\) if no confusion can result. We agree to define the products \(0 \cdot (+\infty)\) and \(0 \cdot (-\infty)\) by giving them the value 0; thus, if f is a numerical function defined on E, and A is a subset of E, \(f\varphi_{A}\) denotes the function that coincides with f on A and is equal to 0 on CA. For every point a of a locally compact space, \(\varepsilon_{a}\) denotes the measure defined by placing a unit mass at the point a (Ch. III, §1, No.~3).

The concept of essential upper integral (resp. essentially integrable function), which will be defined in §1, coincides, as we shall see, with the concept of upper integral (resp. integrable function) when the locally compact space T is countable at infinity (GT, I, §9, No.~9).

The reader who is interested only in integration in locally compact spaces that are countable at infinity may therefore omit reading Nos. 1 to 3 of §1; in the rest of the chapter valid statements are obtained, when the spaces envisioned are countable at infinity, on suppressing the words `essential' and `essentially,' and on replacing the symbol \(\mu^{\bullet}\) by \(\mu^{*}\) and the symbol \(\int^{\bullet}\) by \(\int^{*}\) .

In §§1 to 4, the word measure will always mean positive measure; other measures will be explicitly called not necessarily positive measures or complex measures, as the case may be.

